\section{So sánh Git, GitHub và GitLab}

Nhiều bạn mới học thường gọi gộp cả ba cái tên Git, GitHub và GitLab thành một chữ \texttt{git} cho gọn.
Cách gọi này gây nhầm lẫn khi làm đồ án, vì Git là công cụ chạy trên máy cá nhân,
còn GitHub và GitLab là dịch vụ máy chủ dùng để cộng tác nhóm \cite{progit,atlassiangit}.
Chương này làm rõ sự khác nhau giữa ba khái niệm, so sánh chi tiết GitHub với GitLab
trên hơn hai mươi tiêu chí, đối chiếu các tính năng tương đương,
rồi đưa ra khuyến nghị cụ thể cho đồ án môn học ASPX trên .NET Framework 4.8.1 \cite{msnetfx}.

\subsection{Git, GitHub và GitLab khác nhau thế nào}

Git là hệ thống quản lý phiên bản phân tán, nghĩa là mỗi máy giữ đầy đủ lịch sử commit
và làm việc offline hoàn toàn với các lệnh như \verb|git add|, \verb|git commit|,
\verb|git branch| hay \verb|git log| \cite{progit}.
GitHub và GitLab là hai nền tảng trực tuyến xây dựng quanh Git, bổ sung kho chứa tập trung,
quản lý công việc, review code và CI/CD \cite{githubdocs,gitlabdocs}.
Nói ngắn gọn: Git là cái búa miễn phí ai cũng cầm được,
còn GitHub và GitLab là hai xưởng mộc cùng dùng loại búa đó nhưng cách tổ chức khác nhau.
Bảng \ref{tab:c6-khai-niem} tóm tắt sự khác nhau cốt lõi giữa ba khái niệm này.

\begin{table}[H]
\centering
\caption{Khác nhau cốt lõi giữa Git, GitHub và GitLab}
\label{tab:c6-khai-niem}
\begin{tabular}{p{2.4cm}p{3.7cm}p{3.7cm}p{3.7cm}}
\toprule
Tiêu chí & Git & GitHub & GitLab \\
\midrule
Bản chất & Hệ thống quản lý phiên bản phân tán \cite{progit} & Nền tảng hosting và cộng tác dựa trên Git \cite{githubdocs} & Nền tảng DevOps tích hợp dựa trên Git \cite{gitlabdocs} \\
Chạy ở đâu & Máy cá nhân của từng thành viên & Đám mây GitHub, hoặc máy chủ riêng bản Enterprise & Đám mây GitLab.com, hoặc máy chủ riêng miễn phí \\
Cần mạng không & Không cần cho commit, nhánh, xem log; chỉ cần khi push và pull & Cần mạng cho hầu hết thao tác cộng tác & Cần mạng cho hầu hết thao tác cộng tác \\
Ví dụ thao tác & \verb|git commit|, \verb|git merge| & Tạo PR, review code, chạy Actions & Tạo MR, approve, chạy pipeline \\
Chi phí & Miễn phí, mã nguồn mở & Gói Free, Team, Enterprise & Gói Free, Premium, Ultimate, bản tự quản miễn phí \\
\bottomrule
\end{tabular}
\end{table}

Đoạn lệnh \ref{lst:c6-offline-remote} cho thấy Git local chạy offline hoàn toàn,
chỉ khi đẩy code lên GitHub hoặc GitLab mới cần kết nối mạng.

\begin{lstlisting}[language=bash,caption={Git local lam viec offline, day len remote moi can mang},label={lst:c6-offline-remote}]
# 1. Lam viec offline hoan toan tren may ca nhan
git init AspxDemo_NET481
cd AspxDemo_NET481
git add .
git commit -m "feat: khoi tao project ASPX 4.8.1"
git log --oneline
git branch feature/login
git switch feature/login
# 2. Can mang Internet khi day code len server
git remote add origin https://github.com/nhom3/AspxDemo_NET481.git
git push -u origin feature/login
\end{lstlisting}

Hình \ref{fig:c6-local-remote} minh họa mô hình nhóm 3 người dùng trong đồ án:
mỗi máy chạy Git local, cùng đẩy code lên hai remote là GitHub và GitLab.

\begin{figure}[H]
\centering
\begin{tikzpicture}[>=Stealth, thick, box/.style={draw, rounded corners, fill=blue!10, minimum width=2.6cm, minimum height=0.9cm, align=center, font=\small}]
\node[box] (local) {Git local\\may sinh vien};
\node[box, above right=0.8cm and 1.6cm of local] (hub) {GitHub\\remote chinh};
\node[box, below right=0.8cm and 1.6cm of local] (lab) {GitLab\\remote mirror};
\draw[->] (local) -- node[midway, above, sloped, font=\footnotesize] {push/pull} (hub);
\draw[->] (local) -- node[midway, below, sloped, font=\footnotesize] {push/mirror} (lab);
\end{tikzpicture}
\caption{Mỗi máy chạy Git local và đẩy code lên GitHub và GitLab}
\label{fig:c6-local-remote}
\end{figure}

Điểm mấu chốt cần nhớ: thành thạo Git local là nền tảng bắt buộc của mọi thành viên,
còn GitHub hay GitLab chỉ là nơi đặt kho remote và bộ công cụ cộng tác phía trên.
Mất mạng vẫn commit, tạo nhánh và xem lịch sử bình thường, đó chính là sức mạnh của mô hình phân tán.

\subsection{So sánh chi tiết GitHub và GitLab trên hơn hai mươi tiêu chí}

GitHub và GitLab trùng nhau khoảng 80\% chức năng phổ biến như lưu trữ code,
theo dõi công việc và review, nhưng khác nhau rõ ở triết lý:
GitHub mạnh về cộng đồng mã nguồn mở và hệ sinh thái mở rộng,
GitLab mạnh về quy trình DevOps khép kín và khả năng tự quản máy chủ.
Hai bảng dưới đây so sánh 22 tiêu chí, nhóm theo nền tảng, CI/CD,
cộng tác, bảo mật và chi phí để tiện đối chiếu khi chọn công cụ cho đồ án.

\begin{table}[H]
\centering
\caption{So sánh GitHub và GitLab, phần 1: nền tảng, CI/CD và cộng tác}
\label{tab:c6-gh-gl-1}
\begin{tabular}{p{3.2cm}p{5.4cm}p{5.4cm}}
\toprule
Tiêu chí & GitHub & GitLab \\
\midrule
Bản chất nền tảng & Hosting và cộng tác, mạnh cộng đồng mở \cite{githubdocs} & DevOps tích hợp một cửa, mạnh tự quản \cite{gitlabdocs} \\
Cài đặt self-host & Khó: chỉ bản Enterprise Server, tốn phí và cấu hình nặng & Dễ: bản Community miễn phí, cài một gói Omnibus \\
CI/CD tích hợp & GitHub Actions, workflow chạy theo sự kiện \cite{ghactions} & GitLab CI/CD, pipeline theo stages minh bạch \cite{gitlabci} \\
Runner Windows & Có sẵn image windows-latest, dùng ngay & Tự đăng ký Windows runner, gắn nhãn windows \\
Build ASPX 4.8.1 & Gọi msbuild trên windows-latest, cache NuGet \cite{msmsbuild} & Gọi msbuild trên Windows runner, cache tương tự \\
Review code & Pull request, suggestion, review theo từng file & Merge request, suggestion, review đa cấp \\
Approval bắt buộc & Rulesets và required reviewers, bản trả phí đầy đủ hơn & Approval rules miễn phí đã khá mạnh \\
Bảo vệ nhánh main & Protected branch, status check, CODEOWNERS & Protected branch, protected tag, push rules \\
Quản lý công việc & Issues và Projects dạng kanban & Issues, Boards, Epics và Milestones \\
Wiki và tài liệu & Wiki, Pages và Discussions & Wiki, Pages và Snippets \\
Package NuGet & GitHub Packages hỗ trợ NuGet, npm, Maven & Package Registry hỗ trợ NuGet và Generic \\
\bottomrule
\end{tabular}
\end{table}

\begin{table}[H]
\centering
\caption{So sánh GitHub và GitLab, phần 2: bảo mật, chi phí và phù hợp ASPX}
\label{tab:c6-gh-gl-2}
\begin{tabular}{p{3.2cm}p{5.4cm}p{5.4cm}}
\toprule
Tiêu chí & GitHub & GitLab \\
\midrule
Container registry & Registry ghcr.io gắn với tài khoản và tổ chức & Registry tích hợp sẵn trong từng project \\
Trang tĩnh demo & GitHub Pages, deploy từ nhánh hoặc Actions & GitLab Pages, deploy qua job trong pipeline \\
Quét secret lộ lọt & Secret scanning và push protection & Secret detection chạy trong pipeline \\
Quét bảo mật code & Dependabot và CodeQL \cite{githubdocs} & SAST, DAST và Dependency scanning \cite{gitlabdocs} \\
Trợ lý AI & Copilot gợi ý code trực tiếp trong IDE & Duo gợi ý code và tóm tắt MR \\
Kho mở rộng & Marketplace rất lớn, hàng nghìn Actions sẵn & Tích hợp sẵn nhiều, chợ mở rộng nhỏ hơn \\
Giới hạn gói free & Khoảng 2000 phút Actions mỗi tháng, 500 MB packages & Khoảng 400 phút CI mỗi tháng, 5 GB registry \\
Giá gói trả phí & Team khoảng 4 USD mỗi người mỗi tháng & Premium khoảng 29 USD mỗi người mỗi tháng \\
Mirror và import & Nhập từ GitLab và Bitbucket, mirror hạn chế & Mirror push và pull mạnh, nhập từ GitHub dễ \\
Cộng đồng, việc làm & Rất phổ biến, nhà tuyển dụng hay xem profile & Phổ biến trong doanh nghiệp tự quản máy chủ \\
Phù hợp đồ án ASPX & Rất tốt: Actions Windows sẵn, dễ demo pipeline xanh & Rất tốt: pipeline minh bạch, hợp học DevOps \\
\bottomrule
\end{tabular}
\end{table}

Con số giới hạn và giá ở hai bảng trên thay đổi theo chính sách hãng\footnote{Cần kiểm tra bảng giá chính thức trước ngày bảo vệ: \url{https://github.com/pricing} và \url{https://about.gitlab.com/pricing}.},
nhưng thứ tự lớn nhỏ tương đối giữ nguyên trong nhiều năm:
GitHub hào phóng phút CI miễn phí hơn, GitLab hào phóng tính năng DevOps miễn phí hơn.
Với project ASPX 4.8.1 viết bằng C\# và Web Forms, cả hai nền tảng đều build được
vì đều hỗ trợ runner Windows gọi msbuild \cite{msmsbuild}.
Khác biệt thực tế nằm ở trải nghiệm: Actions cấu hình nhanh trong vài phút,
còn pipeline GitLab mô tả stages rõ ràng nên dễ trình bày trước lớp hơn.

\subsection{Các tính năng tương đương giữa hai nền tảng}

Sinh viên đã quen một nền tảng có thể chuyển sang nền tảng còn lại rất nhanh
vì hầu hết tính năng đều có cặp tương đương, chỉ khác tên gọi và chút ít thao tác.
Bảng \ref{tab:c6-tuong-duong} liệt kê 8 cặp tương đương hay dùng nhất trong đồ án,
kèm ghi chú để khi thầy cô hỏi phản biện có thể trả lời ngay.

\begin{table}[H]
\centering
\caption{Tám cặp tính năng tương đương giữa GitHub và GitLab}
\label{tab:c6-tuong-duong}
\begin{tabular}{p{3.2cm}p{5.4cm}p{5.4cm}}
\toprule
Nhu cầu & GitHub & GitLab \\
\midrule
Đề xuất hợp nhất code & Pull request, gọi tắt PR & Merge request, gọi tắt MR \\
Tích hợp CI/CD & Actions với file YAML trong workflows & Pipelines với file \texttt{.gitlab-ci.yml} \\
Soạn code trên web & Codespaces chạy VS Code trên đám mây & Web IDE gọn nhẹ kèm terminal \\
Quản lý công việc & Issues và Projects & Issues và Boards \\
Chứa gói NuGet & GitHub Packages & Package Registry \\
Trang web tĩnh demo & GitHub Pages & GitLab Pages \\
Quét bảo mật tự động & Dependabot và CodeQL & SAST và Dependency scanning \\
Trợ lý AI gợi ý code & Copilot & Duo \\
\bottomrule
\end{tabular}
\end{table}

Điểm khác tên dễ gây lú nhất là PR và MR: Pull request của GitHub
và Merge request của GitLab thực chất cùng một việc là xin hợp nhất nhánh có review,
chỉ khác tên vì lịch sử hai hãng \cite{githubdocs,gitlabdocs}.
Tương tự, Actions và Pipelines cùng mô tả các bước build, test, deploy bằng YAML,
nên đã viết được file YAML cho bên này thì đọc hiểu file YAML bên kia trong vài phút.
Đoạn lệnh \ref{lst:c6-hai-remote} cho thấy cách gắn cả hai remote vào một repo local
để đẩy code song song, đây là bước chuẩn bị cho mô hình mirror ở mục sau.

\begin{lstlisting}[language=bash,caption={Gan ca hai remote GitHub va GitLab vao mot repo},label={lst:c6-hai-remote}]
# Gan 2 remote: GitHub lam chinh, GitLab lam mirror
git remote add github https://github.com/nhom3/AspxDemo_NET481.git
git remote add gitlab https://gitlab.com/nhom3/AspxDemo_NET481.git
git remote -v
# Day mot nhanh len ca hai noi
git push github main
git push gitlab main
# Mirror tu dong: tren GitLab vao Settings -> Repository -> Mirroring,
# dan URL repo GitHub, che do Push, pipeline chay khi co commit moi
\end{lstlisting}

\subsection{Khuyến nghị cho đồ án môn học ASPX}

Không có đáp án đúng duy nhất cho mọi nhóm, vì lựa chọn phụ thuộc vào mục tiêu demo,
thời gian và máy móc sẵn có.
Phần này đưa ra ba khuyến nghị cụ thể: khi nào nên chọn GitHub,
khi nào nên chọn GitLab, và vì sao nhóm 3 người quyết định dùng cả hai qua cơ chế mirror.

\subsubsection{Khi nào nên chọn GitHub}

\begin{itemize}
\item Muốn demo nhanh trong 5 đến 7 phút: tạo repo, bật Actions với image Windows sẵn,
  pipeline build msbuild xanh ngay mà không cần cài runner riêng \cite{ghactions,msmsbuild}.
\item Nhóm có thành viên mới học Git: tài liệu, video hướng dẫn và câu trả lời trên mạng
  cho GitHub nhiều nhất, tra cứu lỗi nhanh hơn hẳn.
\item Muốn khoe profile khi xin thực tập: nhà tuyển dụng quen xem hồ sơ GitHub,
  số PR và mảng contribution xanh là điểm cộng trực quan.
\item Cần trang demo tĩnh gấp: GitHub Pages xuất bản slide hoặc tài liệu nhóm
  chỉ bằng vài cú nhấp chuột từ nhánh chính.
\end{itemize}

\subsubsection{Khi nào nên chọn GitLab}

\begin{itemize}
\item Muốn học DevOps bài bản: pipeline stages, environments và review apps
  thể hiện trực quan tư duy tích hợp và triển khai liên tục \cite{gitlabci}.
\item Cần approval chặt mà không tốn phí: approval rules và protected branch của GitLab
  ở gói miễn phí đã đủ để demo quy trình review nghiêm túc trước lớp.
\item Phòng lab cho phép tự quản máy chủ: bản Community cài trên một máy tính
  của khoa là cả lớp có forge riêng để thực hành mà không phụ thuộc mạng quốc tế.
\item Muốn một nơi làm tất cả: issues, MR, CI/CD, registry và Pages nằm chung một project,
  đỡ phải mở nhiều dịch vụ rời rạc.
\end{itemize}

\subsubsection{Vì sao nhóm dùng cả hai qua cơ chế mirror}

Nhóm 3 người quyết định lấy GitHub làm kho chính vì Actions Windows cấu hình nhanh,
đồng thời mirror sang GitLab để tận dụng pipeline stages minh bạch khi thuyết trình.
Mỗi thành viên làm trên nhánh feature riêng, mở PR bên GitHub để review,
code sau khi merge tự động mirror sang GitLab và kích hoạt pipeline đối chứng.
Cách làm này cho ba lợi ích rõ rệt: thứ nhất, cả nhóm thực hành được cả PR lẫn MR
nên trả lời tốt mọi câu hỏi phản biện; thứ hai, bảng so sánh trong báo cáo có số liệu thật
đo từ hai pipeline chạy cùng một commit; thứ ba, luôn có một bản sao dự phòng
khi một trong hai dịch vụ nghẽn mạng đúng hôm demo.
Đoạn lệnh \ref{lst:c6-gitlab-yml} là khung pipeline GitLab tối thiểu mà nhóm dùng cho project ASPX,
chạy trên Windows runner và gọi msbuild giống hệt job bên GitHub Actions.

\begin{lstlisting}[language=yaml,caption={Khung pipeline GitLab toi thieu build ASPX 4.8.1},label={lst:c6-gitlab-yml}]
# .gitlab-ci.yml skeleton for ASPX 4.8.1 on Windows runner
stages:
  - build
  - test
variables:
  BUILD_CONFIG: 'Release'
build-aspx:
  stage: build
  tags:
    - windows
  script:
    - nuget restore AspxDemo.sln
    - msbuild AspxDemo.sln /p:Configuration=Release
  artifacts:
    paths:
      - AspxDemo/bin/Release/
test-aspx:
  stage: test
  tags:
    - windows
  script:
    - vstest.console.exe AspxDemo.Tests/bin/Release/AspxDemo.Tests.dll
  dependencies:
    - build-aspx
\end{lstlisting}

Quy trình nhánh của nhóm tuân theo GitFlow thu gọn \cite{gitflow}:
nhánh main luôn build xanh, nhánh develop gom tính năng,
mỗi tính năng một nhánh feature ngắn ngày rồi merge qua PR hoặc MR.
Mọi commit tuân thủ Conventional Commits \cite{convcommits} với tiền tố như feat, fix, docs,
nhờ đó lịch sử đọc như nhật ký dự án thay vì mớ ghi chú lộn xộn.
Tóm lại, GitHub phù hợp khi cần nhanh và khoe được thành quả,
GitLab phù hợp khi cần sâu và trình bày được quy trình,
còn dùng cả hai qua mirror là lựa chọn tối ưu cho bài trình bày so sánh này.
